\documentclass[a4paper]{article}     
\usepackage{amsfonts}
\usepackage{amsmath}
\usepackage{amssymb}
\usepackage{graphicx}

\begin{document}

\noindent \large Auteur: Mohamed Shafik \\
\noindent \large Studierichting: Informatica\\
\noindent \large Oefeningenreeks: 4A nr 24.3\\

\medskip

\normalsize

Bereken $\displaystyle\int \dfrac{x-1}{x^2} \ln x dx$\\

Splits de breuk:\\

\begin{eqnarray*}
\displaystyle\int \dfrac{x-1}{x^2} \ln x dx &=& \displaystyle\int \left(\dfrac{1}{x} - \dfrac{1}{x^2}\right) \ln x dx\\
&=& \displaystyle\int \dfrac{\ln x}{x} dx - \displaystyle\int \dfrac{\ln x}{x^2} dx\\
\end{eqnarray*}

Eerste integraal: substitueer $t = \ln x$, dan is $dt = \dfrac{dx}{x}$\\

\begin{eqnarray*}
\displaystyle\int \dfrac{\ln x}{x} dx &=& \displaystyle\int t dt\\
&=& \dfrac{t^2}{2} + c\\
&=& \dfrac{\ln^2 x}{2} + c\\
\end{eqnarray*}

Tweede integraal: parti\"ele integratie\\

Stel\\

$\left\{\begin{array}{l} u = \ln x \\ dv = \dfrac{dx}{x^2} \end{array}\right. \Rightarrow \left\{\begin{array}{l} du = \dfrac{dx}{x} \\ v = -\dfrac{1}{x} \end{array}\right.$\\

dan wordt\\

\begin{eqnarray*}
\displaystyle\int \dfrac{\ln x}{x^2} dx &=& -\dfrac{\ln x}{x} - \displaystyle\int \left(-\dfrac{1}{x}\right) \cdot \dfrac{dx}{x}\\
&=& -\dfrac{\ln x}{x} + \displaystyle\int \dfrac{dx}{x^2}\\
&=& -\dfrac{\ln x}{x} - \dfrac{1}{x} + c\\
\end{eqnarray*}

Samen:\\

\begin{eqnarray*}
I &=& \dfrac{\ln^2 x}{2} - \left(-\dfrac{\ln x}{x} - \dfrac{1}{x}\right) + c\\
&=& \dfrac{\ln^2 x}{2} + \dfrac{\ln x}{x} + \dfrac{1}{x} + c\\
\end{eqnarray*}

\begin{thebibliography}{9}
\end{thebibliography}
\end{document}
